%%=============================================================================
%% Virtuele simulatie van de straalverbinding (bij ontstentenis van fysieke hardware)
%% Dit fragment is bedoeld voor Hoofdstuk~\ref{ch:proof-of-concept}, als vervolg op
%% Sectie~\ref{sec:plan-van-aanpak-virtualbox} uit de requirementsanalyse.
%%=============================================================================

\section{Volledige virtuele simulatie van de 5.6~GHz airMAX-straalverbinding}%
\label{sec:virtuele-simulatie-airmax}

\subsection{Probleemstelling en methodologische verantwoording}%
\label{subsec:sim-verantwoording}

In Fase~3 van de methodologie (Sectie~\ref{sec:meth-fase3}) werd expliciet gekozen voor
een fysieke proof-of-concept boven een zuiver gesimuleerde netwerkomgeving, net omdat
reële atmosferische demping en hardware-overhead moeilijk betrouwbaar te modelleren
zijn. In de praktijk bleek de aankoop van de vooropgestelde Ubiquiti airMAX-apparatuur
(werkzaam in de 5.6~GHz Wi-Fi-band, zie de Tussentijdse Conclusie in
Sectie~\ref{sec:req-tussentijdse-conclusie}) echter niet haalbaar binnen het tijdsbestek
en budget van deze bachelorproef. Om toch een tastbare, functionele en meetbare
praktijkuitwerking te kunnen voorleggen, wordt de straalverbinding zelf vervangen door
een \emph{volledige netwerkemulatie}: twee virtuele machines die via
\texttt{VirtualBox} gekoppeld worden op een manier die de eigenschappen van de
antenneverbinding zo getrouw mogelijk nabootst, zonder dat er effectief radiogolven aan
te pas komen.

Dit is een erkende en gangbare onderzoeksmethode in netwerkonderzoek: waar de fysieke
laag (modulatie, propagatie, antennegain) niet zelf kan worden opgebouwd, wordt het
\emph{gedrag} van die fysieke laag op de hogere lagen (doorvoersnelheid, latency,
jitter, packet loss) nagebootst met een netwerkemulator, zodat de protocollaag (in dit
geval SFTP over TCP/IP) wél op een realistische manier getest kan worden. Dit is
expliciet \textbf{geen vervanging} van de oorspronkelijke onderzoeksdoelstelling om een
fysieke radioverbinding te realiseren, maar een pragmatische aanpassing van de scope,
die als dusdanig moet worden opgenomen als beperking van dit onderzoek: de resultaten
in Hoofdstuk~\ref{ch:resultaten-en-evaluatie} valideren het functionele en
protocolgedrag van de architectuur onder \emph{verwachte, uit datasheets afgeleide}
kanaalcondities, maar bevestigen niet de effectieve RF-prestaties (antennealignment,
werkelijke propagatie, weersinvloed) van een fysieke airMAX-link in het veld. Dit wordt
expliciet aangehaald als aanbeveling voor toekomstig onderzoek in
Hoofdstuk~\ref{ch:conclusie}.

\subsection{Architectuur: een link-emulator tussen twee eindpunten}%
\label{subsec:sim-architectuur}

De eenvoudigste opzet uit Sectie~\ref{sec:plan-van-aanpak-virtualbox} past
\texttt{tc/netem} rechtstreeks toe op de netwerkinterface van de fileserver-VM. Om de
antenneverbinding echter zo zuiver mogelijk te isoleren van de eindpunten zelf -- net
zoals twee airMAX-radio's elk hun eigen apparaat zijn, met de lucht als gedeeld medium
ertussen -- verdient een \emph{dumbbell-topologie} met een derde, tussenliggende VM de
voorkeur. Figuur~\ref{fig:dumbbell-linkemulator} toont deze opzet.

\begin{figure}[h]
  \centering
  \includegraphics[width=\textwidth]{figs/dumbbell_linkemulator.pdf}
  \caption{Dumbbell-topologie met een aparte link-emulator-VM die de rol van de twee
  airMAX-antennes overneemt. De server- en client-VM draaien onveranderde SFTP-software;
  enkel de link-emulator past \texttt{tc/netem} toe.}
  \label{fig:dumbbell-linkemulator}
\end{figure}

Deze opzet vereist drie VirtualBox-VM's:

\begin{itemize}
  \item \textbf{VM 'server'}: draait de SFTP-server (conform FR-01--FR-03), met één
        netwerkinterface gekoppeld aan \textit{Internal Network A}.
  \item \textbf{VM 'link-emulator'}: een minimale Linux-VM met \emph{twee}
        netwerkinterfaces -- één op \textit{Internal Network A} (richting server), één
        op \textit{Internal Network B} (richting client) -- die als Linux-bridge
        (\texttt{brctl}/\texttt{bridge-utils} of \texttt{systemd-networkd}) fungeert.
        Op beide interfaces wordt onafhankelijk een \texttt{netem}-qdisc toegepast, wat
        toelaat om de verbinding zelfs asymmetrisch te simuleren (zoals in de praktijk
        vaak het geval is bij een reëel airMAX-PtP-link met licht verschillende
        signaalkwaliteit in beide richtingen).
  \item \textbf{VM 'client'}: draait de SFTP-client en testtools, gekoppeld aan
        \textit{Internal Network B}.
\end{itemize}

Deze scheiding heeft als voordeel dat de server- en client-VM's volledig ongewijzigd
blijven -- zij "zien" enkel een gewoon Ethernet-segment, precies zoals ze een
antenne-interface zouden zien in de fysieke opstelling -- terwijl alle kanaalemulatie
geconcentreerd blijft in één centrale, makkelijk aan te passen VM.

\subsection{Van datasheet naar netem-parameters}%
\label{subsec:sim-parameters}

Om de emulatie inhoudelijk te onderbouwen, worden de \texttt{netem}-parameters niet
willekeurig gekozen, maar afgeleid uit de publieke technische specificaties van het
airMAX-platform zelf. Ubiquiti's TDMA-gebaseerde airMAX-protocol vermijdt de
verborgen-node-problematiek van klassieke CSMA/CA-Wi-Fi en levert daardoor
gedocumenteerde prestatiewaarden: in punt-tot-puntmodus haalt airMAX-apparatuur een
reële TCP/IP-doorvoersnelheid van meer dan 150~Mbit/s onder goede
omstandigheden~\autocite{AirMaxTDMA}. De latency van airMAX-AC-apparatuur is rechtstreeks
gekoppeld aan de ingestelde TDD-framingduur: bij een framing van 5~ms bedraagt de
gemiddelde latency ongeveer 10 tot 15~ms, en bij een framing van 8~ms ongeveer 16 tot
24~ms~\autocite{UISPGPSSync}. Deze waarden vormen de basis voor drie representatieve
simulatiescenario's, samengevat in Figuur~\ref{fig:airmax-scenarios}.

\begin{figure}[h]
  \centering
  \includegraphics[width=\textwidth]{figs/airmax_scenarios.pdf}
  \caption{Drie gesimuleerde kanaalscenario's voor de airMAX-link, met
  \texttt{netem}-doelwaarden afgeleid uit de airMAX-datasheets.}
  \label{fig:airmax-scenarios}
\end{figure}

\paragraph{Scenario 1 -- Ideaal (korte afstand, sterk signaal).} Representatief voor het
optimale bereik uit Sectie~\ref{sec:req-tussentijdse-conclusie} (0--15~km, vrije
Fresnelzone): hoge modulatie, 5~ms TDD-framing, verwaarloosbaar packet loss.
\begin{verbatim}
sudo tc qdisc add dev eth1 root netem \
    delay 12ms 2ms distribution normal \
    loss 0.1% \
    rate 150mbit
\end{verbatim}

\paragraph{Scenario 2 -- Marginaal (langere afstand, lagere modulatie).}
Representatief voor het horizonlimiet-segment (15--30~km): de adaptieve modulatie van
airMAX schakelt terug naar een robuustere, minder efficiënte modulatiewijze, de
TDD-framing wordt verlengd, en er treedt meetbaar packet loss op door een lagere
signaal-ruisverhouding.
\begin{verbatim}
sudo tc qdisc add dev eth1 root netem \
    delay 22ms 3ms distribution normal \
    loss 1.5% \
    rate 35mbit
\end{verbatim}

\paragraph{Scenario 3 -- DFS-kanaalwissel (radardetectie).} Zoals toegelicht in
Sectie~\ref{subsec:dfs}, verplicht de detectie van een radarsignaal (bijvoorbeeld van de
weerradar bij Zaventem) een onmiddellijke kanaalwissel. Volgens de Europese
DFS-regelgeving (EN~301~893) moet de zender de uitzending binnen enkele honderden
milliseconden staken (\textit{Channel Move Time}) en mag pas op het nieuwe kanaal
worden uitgezonden na een \textit{Channel Availability Check}, die op de meeste kanalen
zo'n 60~seconden duurt maar op bepaalde radar-gevoelige kanalen kan oplopen tot
10~minuten. Voor de TCP-sessie van SFTP gedraagt dit zich niet als verhoogde latency,
maar als een \emph{volledige, tijdelijke onderbreking} van de verbinding. Dit wordt
gesimuleerd door de \texttt{netem}-qdisc tijdelijk te vervangen door een volledig
pakketverlies, gevolgd door herstel:
\begin{verbatim}
# Simuleer de onderbreking (bv. 10 seconden channel move + hertest)
sudo tc qdisc change dev eth1 root netem loss 100%
sleep 10
sudo tc qdisc change dev eth1 root netem \
    delay 22ms 3ms distribution normal loss 1.5% rate 35mbit
\end{verbatim}
Dit scenario laat toe om FR-06 (hervatten van onderbroken overdrachten) concreet te
testen: een SFTP-sessie die deze onderbreking overleeft en de bestandsoverdracht
nadien hervat, toont aan dat de architectuur bestand is tegen een reële DFS-gebeurtenis.

\subsection{Uitvoering en interpretatie van de resultaten}%
\label{subsec:sim-uitvoering}

Voor elk van de drie scenario's wordt, net zoals in het stappenplan van
Sectie~\ref{sec:plan-van-aanpak-virtualbox}, een representatieve bestandsoverdracht via
SFTP uitgevoerd en gemeten met \texttt{iperf3} (doorvoersnelheid), \texttt{ping}/
\texttt{mtr} (latency) en \texttt{Wireshark} (TCP-venstergedrag, hertransmissies). De
resultaten worden getoetst aan REQ-01 tot en met REQ-03 en aan
Tabel~\ref{tab:requirementsmatrix}, op exact dezelfde manier alsof het om metingen op
een fysieke straalverbinding zou gaan. Waar in Hoofdstuk~\ref{ch:resultaten-en-evaluatie}
naar "de proof-of-concept" wordt verwezen, betreft dit dus deze volledig virtuele
dumbbell-opstelling, met de expliciete kanttekening uit
Sectie~\ref{subsec:sim-verantwoording} dat dit een emulatie van de verwachte
airMAX-prestaties is, en geen fysieke RF-meting.
